\documentclass[10pt,a4paper]{amsart}
\usepackage{amsmath,amssymb,amsthm,graphicx}
% Retained minimal preamble for the Stacks Project chapter varieties.tex.
% Verbatim excerpt of the chapter preamble: Xy-pic setup (native preamble lines 14-18), the native theorem
% declarations with the mathematical macros (native preamble lines 154-217), and the project own no-op
% definitions of the reference/slogan/history display environments (native wrapper.tex lines 1-3) that the
% counted statement uses for its slogan. Copied mechanically; nothing added, removed or redefined.
% For commutative diagrams we use Xy-pic
\usepackage[all]{xy}

% We use 2cell for 2-commutative diagrams.
\xyoption{2cell}
% Theorem environments.
%
\theoremstyle{plain}
\newtheorem{theorem}[subsection]{Theorem}
\newtheorem{proposition}[subsection]{Proposition}
\newtheorem{lemma}[subsection]{Lemma}

\theoremstyle{definition}
\newtheorem{definition}[subsection]{Definition}
\newtheorem{example}[subsection]{Example}
\newtheorem{exercise}[subsection]{Exercise}
\newtheorem{situation}[subsection]{Situation}

\theoremstyle{remark}
\newtheorem{remark}[subsection]{Remark}
\newtheorem{remarks}[subsection]{Remarks}

\numberwithin{equation}{subsection}

% Macros
%
\def\lim{\mathop{\mathrm{lim}}\nolimits}
\def\colim{\mathop{\mathrm{colim}}\nolimits}
\def\Spec{\mathop{\mathrm{Spec}}}
\def\Hom{\mathop{\mathrm{Hom}}\nolimits}
\def\Ext{\mathop{\mathrm{Ext}}\nolimits}
\def\SheafHom{\mathop{\mathcal{H}\!\mathit{om}}\nolimits}
\def\SheafExt{\mathop{\mathcal{E}\!\mathit{xt}}\nolimits}
\def\Sch{\mathit{Sch}}
\def\Mor{\mathop{\mathrm{Mor}}\nolimits}
\def\Ob{\mathop{\mathrm{Ob}}\nolimits}
\def\Sh{\mathop{\mathit{Sh}}\nolimits}
\def\NL{\mathop{N\!L}\nolimits}
\def\CH{\mathop{\mathrm{CH}}\nolimits}
\def\proetale{{pro\text{-}\acute{e}tale}}
\def\etale{{\acute{e}tale}}
\def\QCoh{\mathit{QCoh}}
\def\Ker{\mathop{\mathrm{Ker}}}
\def\Im{\mathop{\mathrm{Im}}}
\def\Coker{\mathop{\mathrm{Coker}}}
\def\Coim{\mathop{\mathrm{Coim}}}

% Boxtimes
%
\DeclareMathSymbol{\boxtimes}{\mathbin}{AMSa}{"02}

%
% Macros for moduli stacks/spaces
%
\def\QCohstack{\mathcal{QC}\!\mathit{oh}}
\def\Cohstack{\mathcal{C}\!\mathit{oh}}
\def\Spacesstack{\mathcal{S}\!\mathit{paces}}
\def\Quotfunctor{\mathrm{Quot}}
\def\Hilbfunctor{\mathrm{Hilb}}
\def\Curvesstack{\mathcal{C}\!\mathit{urves}}
\def\Polarizedstack{\mathcal{P}\!\mathit{olarized}}
\def\Complexesstack{\mathcal{C}\!\mathit{omplexes}}
% \Pic is the operator that assigns to X its picard group, usage \Pic(X)
% \Picardstack_{X/B} denotes the Picard stack of X over B
% \Picardfunctor_{X/B} denotes the Picard functor of X over B
\def\Pic{\mathop{\mathrm{Pic}}\nolimits}
\def\Picardstack{\mathcal{P}\!\mathit{ic}}
\def\Picardfunctor{\mathrm{Pic}}
\def\Deformationcategory{\mathcal{D}\!\mathit{ef}}
\newenvironment{reference}{}{}
\newenvironment{slogan}{}{}
\newenvironment{history}{}{}
\usepackage{cleveref}
\crefname{originallemma}{original lemma}{original lemmas}
\Crefname{originallemma}{Original lemma}{Original lemmas}
\makeatletter
\newlabel{lemma-map-tangent-spaces}{{0B2F}{source}{}{sourcepointer-1}{}}
\newlabel{lemma-map-tangent-spaces@cref}{{[originallemma][0][]0B2F}{[1][0][]source}{}{}{}}
\makeatother
\crefname{originallemma}{original lemma}{original lemmas}
\Crefname{originallemma}{Original lemma}{Original lemmas}
\makeatletter
\newlabel{morphisms-lemma-permanence-finite-type}{{01T8}{source}{}{sourcepointer-2}{}}
\newlabel{morphisms-lemma-permanence-finite-type@cref}{{[originallemma][0][]01T8}{[1][0][]source}{}{}{}}
\makeatother
\crefname{originallemma}{original lemma}{original lemmas}
\Crefname{originallemma}{Original lemma}{Original lemmas}
\makeatletter
\newlabel{morphisms-lemma-triangle-differentials}{{01UX}{source}{}{sourcepointer-3}{}}
\newlabel{morphisms-lemma-triangle-differentials@cref}{{[originallemma][0][]01UX}{[1][0][]source}{}{}{}}
\makeatother
\crefname{originallemma}{original lemma}{original lemmas}
\Crefname{originallemma}{Original lemma}{Original lemmas}
\makeatletter
\newlabel{morphisms-lemma-unramified-at-point}{{02GF}{source}{}{sourcepointer-4}{}}
\newlabel{morphisms-lemma-unramified-at-point@cref}{{[originallemma][0][]02GF}{[1][0][]source}{}{}{}}
\makeatother
\title{Injectivity of the differential on tangent spaces is equivalent to unramifiedness at a point with equal residue fields}
\author{The Stacks Project Authors}
\date{stacks-project:varieties.tex:local-snapshot-20261011; GNU FDL 1.2 or later}
\begin{document}
\maketitle
\noindent\textit{Editorial note.} All delivered mathematical text is the author's own, copied from the pinned Stacks Project chapter file varieties.tex; no statement or proof step was rewritten, reordered or removed.\par\medskip
\noindent\textit{Editorial note.} This card was re-admitted after a gate audit: a text-similarity screen had flagged it against a statement of another chapter, but the two are different propositions, since that one characterizes unramifiedness of a morphism of algebraic spaces through the vanishing of the sheaf of differentials and formal unramifiedness, while this one is a tangent-space criterion at a point. The ruling recorded for this lane is therefore to keep this card.\par\medskip
\noindent\textit{Editorial note.} The counted spans carry 4 navigation stubs and no bibliography citation.\par\medskip
\noindent\textit{Editorial note.} The preamble is the chapter's retained minimal preamble (sources/stacks-varieties/retained-preamble.tex), which since 2026-10-12 also carries the project's own reference/slogan/history definitions.\par\medskip
\noindent\textit{Editorial note.} The label lemma-injective-tangent-spaces-unramified is defined only in varieties.tex, so no same-name restatement exists in another staged chapter.\par\medskip
\noindent\textit{Allowed standard prerequisites (editorial).}\begin{itemize}
\item Duality between the tangent space and the fibre of the cotangent module. Same chapter, Lemma lemma-map-tangent-spaces, invoked at proof lines 3162-3164. Supplies the translation from injectivity of the differential to surjectivity on the fibres of the differentials.
\item Transitivity exact sequence for the module of differentials. Morphisms, Lemma morphisms-lemma-triangle-differentials, invoked at proof lines 3165-3166. Supplies the identification with the vanishing of the fibre of the relative differentials of f.
\item Pointwise criterion for unramifiedness by vanishing of the fibre of the relative differentials. Morphisms, Lemma morphisms-lemma-unramified-at-point, invoked at proof lines 3169-3170. Supplies the last equivalence.
\item Locally finite type is preserved by the morphism. Morphisms, Lemma morphisms-lemma-permanence-finite-type, invoked at proof lines 3160-3161. Supplies that the transitivity sequence applies.
\end{itemize}\medskip
\ifdefined\hypertarget\hypertarget{sourcepointer-1}{}\fi
\noindent\textit{Source reference.} Original lemma 0B2F: sources/stacks-varieties/pinned-original-source.tex, source lines 3108--3108. Pointer outside the counted span, invoked once for context; the load-bearing argument of the counted statement is the author's proof reproduced in full in the proof span.\par\medskip
\ifdefined\hypertarget\hypertarget{sourcepointer-2}{}\fi
\noindent\textit{Source reference.} Original lemma 01T8: sources/stacks-chow/tags.txt, source lines 2365--2365. Pointer outside the counted span, invoked once for context; the load-bearing argument of the counted statement is the author's proof reproduced in full in the proof span.\par\medskip
\ifdefined\hypertarget\hypertarget{sourcepointer-3}{}\fi
\noindent\textit{Source reference.} Original lemma 01UX: sources/stacks-chow/tags.txt, source lines 2426--2426. Pointer outside the counted span, invoked once for context; the load-bearing argument of the counted statement is the author's proof reproduced in full in the proof span.\par\medskip
\ifdefined\hypertarget\hypertarget{sourcepointer-4}{}\fi
\noindent\textit{Source reference.} Original lemma 02GF: sources/stacks-chow/tags.txt, source lines 3226--3226. Pointer outside the counted span, invoked once for context; the load-bearing argument of the counted statement is the author's proof reproduced in full in the proof span.\par\medskip
\begin{lemma}
\label{lemma-injective-tangent-spaces-unramified}
Let $f : X \to Y$ be a morphism of schemes locally of finite type over a
base scheme $S$. Let $x \in X$ be a point. Set $y = f(x)$ and assume
that $\kappa(y) = \kappa(x)$. Then the following are equivalent
\begin{enumerate}
\item $\text{d}f : T_{X/S, x} \longrightarrow T_{Y/S, y}$ is injective, and
\item $f$ is unramified at $x$.
\end{enumerate}
\end{lemma}

\begin{proof}
The morphism $f$ is locally of finite type by
Morphisms, Lemma \ref{morphisms-lemma-permanence-finite-type}.
The map $\text{d}f$ is injective, if and only if
$\Omega_{Y/S, y} \otimes \kappa(y) \to \Omega_{X/S, x} \otimes \kappa(x)$
is surjective (Lemma \ref{lemma-map-tangent-spaces}).
The exact sequence $f^*\Omega_{Y/S} \to \Omega_{X/S} \to \Omega_{X/Y} \to 0$
(Morphisms, Lemma \ref{morphisms-lemma-triangle-differentials})
then shows that this happens if and only if
$\Omega_{X/Y, x} \otimes \kappa(x) = 0$.
Hence the result follows from
Morphisms, Lemma \ref{morphisms-lemma-unramified-at-point}.
\end{proof}

\end{document}
